\section{결론}
\label{sec:conclusion}
우리는 군중 내비게이션 로봇이 할 수 있는 것 -- 사람 보행 행동의 코드북에서 디코딩한 안전한 경로 -- 과, 그중 특정 사용자가
원하는 것을 분리하고, 후자를 소수의 A/B 질문으로부터 학습하였다. 다섯 번의 EUBO 질문은 처음 보는 상황에서의 단순 regret을
인간 데이터 사전분포의 수준(0.45--0.47)에서 0.09--0.12로 줄인다. EUBO든 정보 이득이든 사후분포로 질문을 선택하는 것이 추정기보다 더 중요하며, 후보 집합은
어떤 선호로도 도달할 수 있는 수준을 제한한다. 다음으로, Sec.~\ref{sec:pilot}의 사용자 연구는 실제 사람들이 충분히 달라서
개인 $\hat\theta$가 공유된 것보다 학습에 쓰지 않은 답변을 더 잘 예측하는지를 검증할 것이며, 반응하는 보행자와 학습된
예측기를 사용한 폐루프 테스트는 군중이 로봇에 반응할 때 선호가 무엇을 바꾸는지를 보여줄 것이다.
